% GERADO por tools/generation/gerar_secoes_latex.py a partir de
% artifacts/paper/secoes/02-referencial.md. Nao editar a mao: a proxima geracao
% desfaz. Corrija o Markdown e regere.
\subsection{Observatórios de projetos}

Não há na literatura definição única para \textit{observatório}. Os estudos convergem em associá-lo a ``instrumentos voltados à observação e à promoção da transparência'', com atividades de coleta, consolidação, armazenamento, análise e divulgação de dados sobre um setor ou área do conhecimento \cite{sakata2013construccao,tinati2015building,keever2017decada}, e um observatório se materializa como ``unidade organizacional, quando o observatório se configura como um elemento formal da organização'', como processo pelo qual um grupo executa a observação, ou como instrumento tecnológico \cite{vieira2020universal}. Observatórios de projetos se distinguem pelo objeto: têm como finalidade ``a disponibilização sistemática, objetiva e detalhada de informações específicas sobre projetos'' \cite{vieira2021model}.

O termo é de nicho no campo de projetos, e convém demarcá-lo frente aos vizinhos com que se confunde. O \textbf{gerenciamento de projetos} exerce autoridade sobre o projeto --- planeja, decide e age sobre escopo, prazo, custo e riscos, respondendo por seus resultados \cite{pmi2021pmbok,turner2016gower} ---, e um observatório não gerencia os projetos que observa nem responde por eles. O \textbf{monitoramento e controle de projetos} compara o desempenho observado à linha de base para provocar ação corretiva, e sua audiência é a equipe e o patrocinador; o observatório não fecha essa malha --- seu produto é informação divulgada, sua audiência é externa à execução e seu escopo é um conjunto de projetos, não um só, e ferramentas de monitoramento governamental servem ao gestor que executa \cite{trois2017transparency}, não a quem o fiscaliza. O \textbf{escritório de projetos (PMO)} padroniza, apoia ou dirige a governança de projetos \cite{pmi2021pmbok}, autoridade de que o observatório carece --- mas as figuras são compatíveis: um PMO pode manter um observatório, e um observatório pode reunir projetos de várias organizações, sem PMO algum. A categoria é geral --- há observatórios urbanos \cite{keever2017decada}, turísticos \cite{pimentel2018observatorio} e de saúde \cite{yoshiura2018towards} ---, e o recorte aqui é o objeto.

\subsection{O Modelo para Observatórios de Projetos}

O MPO é um modelo conceitual de referência para a concepção e a análise de observatórios de projetos, publicado em versão preliminar \cite{vieira2021model}, avaliado e evoluído em estudos posteriores \cite{vieira2022evaluating,vieira2023survey} e consolidado em sua versão final \cite{vieira2022thesis}. Ele reúne 61 conceitos em três níveis e três dimensões, cada uma dividida em subdimensões, e toda a descrição é textual: cada elemento tem definição em prosa e uma coluna de relações também em prosa, redigida com verbos como \textit{Suporta}, \textit{Viabiliza} e \textit{Especifica}.

A dimensão \textbf{Estruturas} ``reúne os elementos relacionados à configuração estrutural e estática dos observatórios de projetos'' em três subdimensões. \textit{Componentes} agrupa os elementos de Coleta, Processamento e Armazenamento dos dados, o de Disseminação e o de Relacionamento, que coordena a interação entre usuários, projetos e o observatório, mais a infraestrutura de TI que os sustenta --- Redes, Hardware, Serviços, Gerenciamento de dados e Software. \textit{Conteúdos} abrangem os projetos, suas temáticas, os usuários e o próprio observatório; \textit{Características}, as nove propriedades dos dados armazenados, entre elas interoperabilidade, acesso e segurança.

A dimensão \textbf{Processos} descreve os procedimentos executados no observatório: a subdimensão de entrada e saída de dados detalha Coletar, Tratar, Armazenar e Disponibilizar, e a de processamento reúne outros onze, de Transformar a Responsabilizar. Vários nomes reaparecem nas duas dimensões --- Coleta, Processamento e Armazenamento são elementos de Estruturas, e Coletar, Tratar e Armazenar, de Processos ---, cada verbo descrito como realizado por meio do componente correspondente.

A dimensão \textbf{Agentes} ``trata dos participantes envolvidos nos observatórios de projetos, considerando tanto sua natureza quanto os fatores que motivam sua interação com o observatório'': a subdimensão \textit{Atores} enumera as partes interessadas dos projetos, a equipe de gestão e desenvolvimento, os usuários do observatório e os sistemas computacionais externos que trocam dados com ele, e a subdimensão \textit{Motivações} lista onze fatores que os levam a interagir com ele. São esses elementos, e as relações que a prosa enuncia entre eles, que a análise da \S{}5 discute.

\subsection{Representation Theory}

A Representation Theory, de Wand e Weber, é a teoria de Sistemas de Informação que sustenta a análise. Seu ponto de partida é que um sistema de informação é uma representação de um domínio do mundo real, e que a estrutura profunda do sistema determina sua capacidade de informar \cite{wand1995deep}. Daí o critério aplicado a gramáticas de modelagem: uma gramática descreve bem o mundo na medida em que seus construtos correspondem, um a um, aos de uma ontologia de referência \cite{wand1993ontological,weber1997ontological}. Duas correspondências são examinadas --- o \textit{mapeamento de representação}, dos construtos ontológicos aos da gramática, e o \textit{mapeamento de interpretação}, no percurso inverso ---, e, quando alguma deixa de ser total e unívoca, a teoria nomeia quatro deficiências:

\textbf{déficit de construto} (\textit{construct deficit}), quando há construto ontológico que a gramática não expressa; \textbf{sobrecarga de construto} (\textit{construct overload}), quando um construto da gramática representa dois ou mais construtos ontológicos e distinções do domínio colapsam num símbolo só; \textbf{redundância de construto} (\textit{construct redundancy}), quando dois ou mais construtos representam o mesmo construto ontológico; e \textbf{excesso de construto} (\textit{construct excess}), quando um construto da gramática não corresponde a construto ontológico algum. A teoria prevê que se manifestem como ambiguidade, perda de informação e dependência de convenções não documentadas, e há evidência empírica de que afetam o uso efetivo de modelos por seus destinatários \cite{recker2011ontological}. O que a tipologia classifica é o \textit{mapeamento} entre uma gramática e uma ontologia, e usá-la exige declarar qual é o par sob exame --- aqui, o modelo de referência em prosa e a formalização que dele se derivou (\S{}4).

\subsection{UFO e OntoUML}

A Unified Foundational Ontology é uma ontologia de fundamentação que fornece as categorias e as restrições contra as quais modelos de domínio são construídos e avaliados \cite{guizzardi2005ontological,guizzardi2022ufo}, organizada em microteorias. A \textbf{UFO-A} trata dos endurantes e os distingue por \textit{identidade}, \textit{rigidez} e \textit{dependência}; desses critérios saem os sortais («kind», «subkind», «phase», «role»), os não-sortais («category», «mixin», «roleMixin»), os momentos («relator», «mode») e os coletivos («collective»), além das relações de parte-todo, que distinguem a composição entre complexos funcionais («componentOf») da adesão de um membro a um coletivo («memberOf»). A \textbf{UFO-B} trata dos perdurantes: eventos têm partes temporais e contam com a participação de endurantes, o que permite responder \textit{quando} algo ocorreu e \textit{quem} dele participou. A \textbf{UFO-C} trata das entidades sociais e intencionais: distingue agentes de objetos não agentivos pela posse de momentos intencionais e, entre os agentes, os físicos dos sociais.

A OntoUML dá forma a essas categorias como estereótipos de um perfil UML, impondo aos modelos as restrições correspondentes \cite{guizzardi2008grounding,ontouml_readthedocs}. Modelar nela obriga a decidir, para cada classe, qual categoria ontológica ela instancia --- obrigação que converte escolhas implícitas de um texto em compromissos explícitos e verificáveis por ferramenta.
